% \iffalse
%
% hamacro.dtx
% 俺のLaTeXマクロ集(煮ても焼いてもかまいません)
% hamano@cuspy.org
%
%<*driver>
\ProvidesFile{hamacro.dtx}
%</driver>
%<package>\NeedsTeXFormat{LaTeX2e}[1999/12/01]
%<package>\ProvidesPackage{hamacro}[2004/11/05 hamano]
%<*package>
%</package>
%
%<*driver>
%\documentclass{ltxdoc}
\documentclass[a4paper,12pt,uplatex]{jsarticle}
\usepackage{doc}
\usepackage[blue]{hamacro}
\usepackage{layout}
\EnableCrossrefs
\CodelineIndex
\RecordChanges
\setcounter{StandardModuleDepth}{1}

%\changes{v1.0}{2004/11/05}{Initial version}
\GetFileInfo{hamacro.dtx}
\DoNotIndex{\newcommand,\newenvironment}

\begin{document}
  \DocInput{hamacro.dtx}
  \PrintChanges
  \PrintIndex
\end{document}
%</driver>

% fontencを使うと何故かビットマップフォントが埋め込まれるので必ず欧文
% フォントを指定すること
\usepackage[T1]{fontenc}
\usepackage{newtxmath}
%\renewcommand{\rmdefault}{ptm} % Nimbus Roman No9 L
%\renewcommand{\sfdefault}{phv} % Nimbus Sans L
\usepackage{tgtermes}
\usepackage{tgheros}
\renewcommand{\bfdefault}{bx} % 和文がBoldにならない問題回避、tgtermesのバグ? 
\usepackage{inconsolata} % only ttdefault
\usepackage[deluxe]{otf}
\usepackage[dvipdfmx]{graphicx}

% 色定義
\usepackage[dvipdfmx]{color}
\usepackage[dvipdfmx]{xcolor}
\definecolor{lightgray}{HTML}{D3D3D3}
\definecolor{gainsboro}{HTML}{DCDCDC}
\definecolor{azure}{HTML}{F0FFFF}
\definecolor{darkblue}{HTML}{00008B}
\definecolor{midnightblue}{HTML}{191970}
\definecolor{lightskyblue}{HTML}{87CEFA}
\definecolor{olivegreen}{cmyk}{0.64,0,0.95,0.40}
\definecolor{cadetblue}{cmyk}{0.62,0.57,0.23,0}
\definecolor{brown}{cmyk}{0,0.81,1,0.60}

\definecolor{lightgreen}{HTML}{90EE90}
\definecolor{green}{HTML}{008000}
\definecolor{darkgreen}{HTML}{006400}
\definecolor{samuraiblue}{HTML}{1D327D}

\usepackage[dvipdfm]{hyperref}
\hypersetup{
  setpagesize=false,
  bookmarksnumbered=true,
  colorlinks=true,
  linkcolor=linkcolor,
  citecolor=linkcolor,
  filecolor=linkcolor,
  urlcolor=linkcolor
}

% しおりの文字化け対策
\usepackage{pxjahyper}

% その他パッケージ
\usepackage{here}
\usepackage{ctable}
\usepackage{longtable}
\usepackage{ulem}
\usepackage{soul}

% カラーテーマ設定
\colorlet{framerule}{gray}
\colorlet{framebg}{gainsboro}
\colorlet{linkcolor}{blue}

\colorlet{accentcolor}{black}
\colorlet{subcolor}{black}

\DeclareOption{green}{
  \colorlet{accentcolor}{green}
  \colorlet{subcolor}{darkgreen}
}
\DeclareOption{blue}{
  \colorlet{accentcolor}{lightskyblue}
  \colorlet{subcolor}{midnightblue}
}
\DeclareOption{samuraiblue}{
  \colorlet{accentcolor}{lightskyblue}
  \colorlet{subcolor}{samuraiblue}
}
\ProcessOptions

% jsclassesとgeometryパッケージは相性が悪い
%\usepackage[a4paper,top=35mm, bottom=30mm, left=30mm, right=30mm]{geometry}
\setlength{\oddsidemargin}{-1truein}
\addtolength{\oddsidemargin}{0.11\paperwidth}
\setlength{\evensidemargin}{-1truein}
\addtolength{\evensidemargin}{0.11\paperwidth}
\setlength{\marginparwidth}{0.10\paperwidth}
\setlength{\marginparsep}{0.005\paperwidth}
\setlength{\textwidth}{\paperwidth}
\addtolength{\textwidth}{-0.22\paperwidth}

\message{^^J}
\message{***** Layout Info *****^^J}
\message{\hoffset: \the\hoffset^^J}
\message{\voffset: \the\voffset^^J}
\message{\paperwidth: \the\paperwidth^^J}
\message{\paperheight: \the\paperheight^^J}
\message{\textwidth: \the\textwidth^^J}
\message{\textheight: \the\textheight^^J}
\message{\oddsidemargin: \the\oddsidemargin^^J}
\message{\evensidemargin: \the\evensidemargin^^J}
\message{\marginparwidth: \the\marginparwidth^^J}
\message{\marginparsep: \the\marginparsep^^J}
\message{***********************^^J}

% default title formatting
\def\@maketitle{
  \newpage
  \null
  \begin{center}
    {\LARGE\gtfamily\bfseries{\@title}\par}
  \end{center}
  \begin{flushright}
    {\large \@date \par}
    {\large \@author \par}
  \end{flushright}
  \vspace{0.5in}
}

% chapter formatting
\usepackage[Conny]{fncychap}
\ChNameVar{\centering\Huge\sffamily\color{black}}
\ChTitleAsIs
\ChTitleVar{\centering\Huge\sffamily\color{black}}

\renewcommand{\DOCH}{
  \color{subcolor}
  \mghrulefill{3\RW}\par\nobreak
  \vskip -0.5\baselineskip
  \mghrulefill{\RW}\par\nobreak
  \CNV\FmN{\@chapapp}\space \CNoV\thechapter \postchaptername
  \par\nobreak
  \vskip -0.5\baselineskip
}
\renewcommand{\DOTI}[1]{
  \color{subcolor}
  \mghrulefill{\RW}\par\nobreak
  \CTV\FmTi{#1}\par\nobreak
  \vskip 1\baselineskip
}
\renewcommand{\DOTIS}[1]{
  \color{subcolor}
  \mghrulefill{\RW}\par\nobreak
  \CTV\FmTi{#1}\par\nobreak
  \vskip 1\baselineskip
}

% titlesec settings
\usepackage{etex}
\usepackage[explicit,calcwidth]{titlesec}

%% section formatting
\newcommand\sectionformat[2]{
  \tikz {
    \node [
      draw=subcolor,
      line width=1.5pt,
      fill=gray!20,
      align=left,
      rounded corners=6pt,
      text width=\textwidth
    ] {\textcolor{subcolor}{#2}\;\textcolor{black}{#1}};
  }
}

\titleformat{name=\section}
            {\sf\gtfamily\bfseries\LARGE}{}{0pt}
            {\sectionformat{#1}{\textrm\thesection}}
\titleformat{name=\section,numberless}
            {\sf\gtfamily\bfseries\LARGE}{}{0pt}
            {\sectionformat{#1}{}}
% \titlespacing{command}{left}{before-sep}{after-sep}{right-sep}
\titlespacing*{\section}{-1em}{3ex}{-4ex}

% subsection formatting
\titleformat{name=\subsection}[block]
            {\sf\gtfamily\bfseries\Large}
            {\textcolor{subcolor}{\thesubsection}}
            {0.5em}
            {#1}
            [{\addtolength{\titlewidth}{0.2em}
                \titleline*[l]{\color{subcolor}\titlerule[1.5pt]}}]
\titleformat{name=\subsection,numberless}[block]
            {\sf\gtfamily\bfseries\Large}
            {}
            {0.5em}
            {#1}
            [{\addtolength{\titlewidth}{0.2em}
                \titleline*[l]{\color{subcolor}\titlerule[1.5pt]}}]

% subsubsection formatting
\titleformat{\subsubsection}[block]
            {\sf\gtfamily\bfseries\large}
            {\color{subcolor}{\thesubsubsection}}
            {0.5em}
            {#1}
\titleformat{name=\subsubsection,numberless}[block]
            {\sf\gtfamily\bfseries\large}
            {}
            {0.5em}
            {#1}

%% paragraph formatting
\setcounter{secnumdepth}{4}
\titleformat{name=\paragraph}[block]
            {\sf\gtfamily\bfseries\normalsize}{\color{subcolor}{\theparagraph} \color{black}{#1}}{1em}{}
\titleformat{name=\paragraph,numberless}[block]
            {\sf\gtfamily\bfseries\normalsize}{#1}{1em}{}

% 修飾マクロ
\let\origemph\emph
\newcommand{\redemph}[1]
           {\textcolor{red}{\underline{\origemph{\textcolor{black}{#1}}}}}

\newcommand{\ulemph}[1]
           {\textbf{\Underline{#1}}}

\newcommand{\bouemph}[1]
           {\origemph{\bou{#1}}}

%\let\emph\bouemph

% tableofcontentsをサンセリフに統一
\let\origtableofcontents\tableofcontents
\renewcommand{\tableofcontents}{\textsf\origtableofcontents}

% setting for mdframed
\usepackage[framemethod=TikZ]{mdframed}
\mdfsetup{
  middlelinewidth=2pt,
  middlelinecolor=gray,
  backgroundcolor=gainsboro,
  roundcorner=3pt}

% style for quote
%\usepackage{framed}
\usepackage{anyfontsize}
\newcommand{\bigquotemark}{
  {\fontsize{3em}{3em}\selectfont \color{accentcolor} \textbf{``}}
  \vskip -2.3em
  \hskip 1.5zw
}
\renewenvironment{quote}{
  \begin{mdframed}
  \parindent 0zw
  \vskip -0.5em
%  \hskip -1zw
  {\bigquotemark}
}
{
  \smallskip
  \end{mdframed}
}

% style for verbatim

\renewenvironment{verbatim}{
  \VerbatimEnvironment
  \begin{mdframed}
  \tt\fontsize{11.2truept}{1.2em}
  \begin{Verbatim}}
{
  \end{Verbatim}
  \end{mdframed}
}

\newenvironment{verbatim60}{
  \VerbatimEnvironment
  \begin{mdframed}
  \tt\fontsize{15truept}{1.5em}
  \begin{Verbatim}}
{
  \end{Verbatim}
  \end{mdframed}
}

\newenvironment{verbatim40}{
  \VerbatimEnvironment
  \begin{mdframed}
  \tt\fontsize{22truept}{1.5em}
  \begin{Verbatim}}
{
  \end{Verbatim}
  \end{mdframed}
}

% style for pandoc code highlighting
\usepackage{fancyvrb}
\newenvironment{Shaded}{\begin{mdframed}}{\end{mdframed}}
\DefineVerbatimEnvironment{Highlighting}{Verbatim}{commandchars=\\\{\}}
\newcommand{\KeywordTok}[1]{\textcolor[rgb]{0.00,0.44,0.13}{\textbf{{#1}}}}
\newcommand{\DataTypeTok}[1]{\textcolor[rgb]{0.56,0.13,0.00}{{#1}}}
\newcommand{\DecValTok}[1]{\textcolor[rgb]{0.25,0.63,0.44}{{#1}}}
\newcommand{\BaseNTok}[1]{\textcolor[rgb]{0.25,0.63,0.44}{{#1}}}
\newcommand{\FloatTok}[1]{\textcolor[rgb]{0.25,0.63,0.44}{{#1}}}
\newcommand{\ConstantTok}[1]{\textcolor[rgb]{0.53,0.00,0.00}{{#1}}}
\newcommand{\CharTok}[1]{\textcolor[rgb]{0.25,0.44,0.63}{{#1}}}
\newcommand{\SpecialCharTok}[1]{\textcolor[rgb]{0.25,0.44,0.63}{{#1}}}
\newcommand{\StringTok}[1]{\textcolor[rgb]{0.25,0.44,0.63}{{#1}}}
\newcommand{\VerbatimStringTok}[1]{\textcolor[rgb]{0.25,0.44,0.63}{{#1}}}
\newcommand{\SpecialStringTok}[1]{\textcolor[rgb]{0.73,0.40,0.53}{{#1}}}
\newcommand{\ImportTok}[1]{{#1}}
\newcommand{\CommentTok}[1]{\textcolor[rgb]{0.38,0.63,0.69}{\textit{{#1}}}}
\newcommand{\DocumentationTok}[1]{\textcolor[rgb]{0.73,0.13,0.13}{\textit{{#1}}}}
\newcommand{\AnnotationTok}[1]{\textcolor[rgb]{0.38,0.63,0.69}{\textbf{\textit{{#1}}}}}
\newcommand{\CommentVarTok}[1]{\textcolor[rgb]{0.38,0.63,0.69}{\textbf{\textit{{#1}}}}}
\newcommand{\OtherTok}[1]{\textcolor[rgb]{0.00,0.44,0.13}{{#1}}}
\newcommand{\FunctionTok}[1]{\textcolor[rgb]{0.02,0.16,0.49}{{#1}}}
\newcommand{\VariableTok}[1]{\textcolor[rgb]{0.10,0.09,0.49}{{#1}}}
\newcommand{\ControlFlowTok}[1]{\textcolor[rgb]{0.00,0.44,0.13}{\textbf{{#1}}}}
\newcommand{\OperatorTok}[1]{\textcolor[rgb]{0.40,0.40,0.40}{{#1}}}
\newcommand{\BuiltInTok}[1]{{#1}}
\newcommand{\ExtensionTok}[1]{{#1}}
\newcommand{\PreprocessorTok}[1]{\textcolor[rgb]{0.74,0.48,0.00}{{#1}}}
\newcommand{\AttributeTok}[1]{\textcolor[rgb]{0.49,0.56,0.16}{{#1}}}
\newcommand{\RegionMarkerTok}[1]{{#1}}
\newcommand{\InformationTok}[1]{\textcolor[rgb]{0.38,0.63,0.69}{\textbf{\textit{{#1}}}}}
\newcommand{\WarningTok}[1]{\textcolor[rgb]{0.38,0.63,0.69}{\textbf{\textit{{#1}}}}}
\newcommand{\AlertTok}[1]{\textcolor[rgb]{1.00,0.00,0.00}{\textbf{{#1}}}}
\newcommand{\ErrorTok}[1]{\textcolor[rgb]{1.00,0.00,0.00}{\textbf{{#1}}}}
\newcommand{\NormalTok}[1]{{#1}}

% fix for pandoc 1.14
\providecommand{\tightlist}{
  \setlength{\itemsep}{0pt}\setlength{\parskip}{0pt}}

% style for listings
\usepackage{upquote}
\usepackage{listings}
%\usepackage{jlisting}

% ハイフンをU+002D HYPHEN-MINUSのままで表示
%\makeatletter
%\lst@CCPutMacro\lst@ProcessOther {"2D}{\lst@ttfamily{-{}}{-{}}}
%\@empty\z@\@empty
%\makeatother

% アスタリスクをU+002A ASTERISKのままで表示
%\makeatletter
%\lst@CCPutMacro\lst@ProcessOther {"2A}{\lst@ttfamily{*{}}{*{}}}
%\@empty\z@\@empty
%\makeatother

% 一部のPDFリーダーで白い線が入ってしまう問題があるのでmdframeでラップする
% http://tex.stackexchange.com/questions/196247/lstlisting-width-of-background-box

% captionもmdframe内に入るのでいまいち
\AtBeginDocument{
  \def\lstlistingname{リスト}
}

\mdfsetup{
  linewidth=1truept,
  linecolor=gray,
  innerleftmargin=0.5zw,
  innerrightmargin=0.5zw,
  backgroundcolor=gainsboro
}

\BeforeBeginEnvironment{lstlisting}{
  \begin{mdframed}
}

\AfterEndEnvironment{lstlisting}{
  \end{mdframed}
}

\lstset{
  frame=single,
  rulecolor=\color{gainsboro},
  backgroundcolor=\color{gainsboro},
%  numbers=left,
%  numbersep=0pt,
  framesep=0pt,
  xleftmargin=0pt,
  basicstyle=\fontsize{11.2truept}{1.2em}\tt,
  belowskip=-0.5em,
  identifierstyle=\color{black},
  commentstyle=\color{brown},
  keywordstyle=\color{olivegreen},
  ndkeywordstyle=\color{red},
  stringstyle=\color{purple},
  showstringspaces=false,
  breaklines=true,
  columns=[l]{fullflexible}
}

\lstdefinestyle{command}{
  language=command,
  commentstyle=\color{black}
}

\lstdefinelanguage{command}
{
  keywords={\#, export, if, while, for},
  morekeywords={\$},
  otherkeywords={\#},
  morestring=[b]',
  morestring=[b]\"
}

\lstdefinelanguage{JavaScript}{
  keywords={typeof, new, true, false, catch, function, return, null, catch, switch, var, if, in, while, do, else, case, break},
  ndkeywords={class, export, boolean, throw, implements, import, this},
  sensitive=false,
  comment=[l]{//},
  morecomment=[s]{/*}{*/},
  morestring=[b]',
  morestring=[b]\"
}

\lstdefinelanguage{httpd.conf}
{
  morekeywords={VirtualHost, Include, Directory, LogFormat, CustomLog},
  morecomment=[l]{\#}
}

\lstdefinelanguage{nginx.conf}
{
  morekeywords={
    server,server_name,listen,location,
    rewrite,proxy_pass, proxy_set_header},
  morecomment=[l]{\#}
}

\lstdefinelanguage{slapd.conf}
{
  morekeywords={
    include,loglevel,timelimit,sizelimit,
    database,suffix,rootdn,rootpw,directory,
    overlay,index,
  },
  morestring=[b]\",
  morecomment=[l]{\#}
}

% \fi
%
% \iffalse
%<*doc>
% \fi

\title{\LaTeX{} hamacro パッケージ}
\author{hamano}
\maketitle
\tableofcontents
\clearpage
%\MakeShortVerb{\|}

\section{はじめに}

これは私が本や文書を組版する際に利用する\LaTeX{}のスタイルパッケージで
す。

80文字ピッタリおさまるverbatim環境や、各種プログラミング言語のコードハ
イライティングなど、主にプログラミングに関する文書に適しています。

既存の\LaTeX{}文書で簡単に利用できることと、pandocやRe:Viewが出力する
\LaTeX{}で利用することを想定していますので、新しいマクロを定義するとい
うよりも既存のマクロや環境を書き換えています。

このパッケージはMITライセンスで配布していますので、利用、改変、部分的な
コピーなど自由です!

\section{使い方}

hamacro.styを然るべき場所に配置し、プリアンブルに

\begin{lstlisting}[language=tex]
\usepackage{hamacro}
\end{lstlisting}

と記述するだけでこのドキュメントの様なスタイルで組版できます。

具体的な用例は\href{https://github.com/hamano/hamacro/blob/master/example.tex}{example.tex}を参照して下さい。

\subsection{pandocでの利用方法}

このスタイルはpandocで生成したTeX文書に最適化しています。
ですのでmarkdownで書いた文書を手軽にPDFに組版できます。
例えば\href{https://github.com/hamano/hamacro/blob/master/examples/memo/memo.pdf}{memo.md}を変換して\href{https://github.com/hamano/hamacro/blob/master/examples/memo/memo.pdf}{memo.pdf}を生成できます。

詳しくはこちらのサンプルディレクトリを見てください。

 - \href{https://github.com/hamano/hamacro/blob/master/examples/memo/}{https://github.com/hamano/hamacro/blob/master/examples/memo/}

\subsection{Re:Viewでの利用方法}

\documentclass[a4paper,onecolumn,12pt]{jsarticle}

\usepackage{hamacro}
\hypersetup{
  pdfstartview={FitH},
  pdftitle={タイトル},
  pdfauthor={著者},
  pdfsubject={サブタイトル},
  pdfkeywords={キーワード},
}

% title
\title{\LaTeX テンプレート}
\author{HAMANO Tsukasa}
\date{\today}

% document
\begin{document}
\maketitle
\vspace{1em}
\tableofcontents
\clearpage

\section{修飾}

formatting by \emph{emph}

formatting by \textbf{textbf}

emphによる\emph{強調}

textbfによる\textbf{強調}

\subsection{フォントサイズ}

\tiny{あ tiny}\\
\scriptsize{あ scriptsize}\\
\footnotesize{あ footnotesize}\\
\small{あ small}\\
\normalsize{あ normalsize}\\
\large{あ large}\\
\Large{あ Large}\\
\LARGE{あ LARGE}\\
\huge{あ huge}\\
\Huge{あ Huge}\\
\normalsize{リセット}\\

\subsection{リンク}
\url{http://localhost/}

\href{http://www.example.com/}{example.com}

\section{リスト}

\subsection{itemaize}
\begin{itemize}
\item その１
\item その２
\item その３
\end{itemize}

\subsection{enumerate}
\begin{enumerate}
\item その１
\item その２
\item その３
\end{enumerate}

\subsection{description}
\begin{description}
  \item[名前１] 説明１
  \item[名前２] 説明２
  \item[名前３] 説明３
\end{description}

\section{引用}

\subsection{quote}
字下げなし
\begin{quote}
１行目

２行目

３行目
\end{quote}

\subsection{quotation}
字下げあり
\begin{quotation}
１行目

２行目

３行目
\end{quotation}

\section{ソースコード}


hello.c

\begin{verbatim}
#include <stdio.h>

int main(int argc, char *argv[]){
    printf("Hello World!\n");
    return 0;
}
\end{verbatim}

\section{テーブル}
\subsection{普通のテーブル}
\begin{tabular}{|l|c|r|}
\hline
セル1 & セル2 & セル3 \\ \hline
セル4 & セル5 & セル6 \\ \hline
セル7 & セル8 & セル9 \\ \hline
\end{tabular}

\subsection{タイトル付きテーブル(中央)}

\begin{table}[htbp]
\begin{center}
\begin{tabular}{|l|c|r|}
\hline
セル1 & セル2 & セル3 \\ \hline
セル4 & セル5 & セル6 \\ \hline
セル7 & セル8 & セル9 \\ \hline
\end{tabular}
\caption{タイトル}
\end{center}
\end{table}

\subsection{水平線}
\begin{center}\rule{3in}{0.4pt}\end{center}

\clearpage
\section{数式}
\begin{enumerate}

\item 質量とエネルギーの等価性
\begin{equation}
E = mc ^2
\end{equation}

\item 波動方程式
\begin{equation}
\frac{\partial^2 z}{\partial t^2}=
c^2 (\frac{\partial^2 z}{\partial x^2}+\frac{\partial^2 z}{\partial y^2})-
\mu \frac{\partial z}{\partial t}
\end{equation}

\begin{equation}
\sum_{x=1}^4 x^2 = 30
\end{equation}

\item オイラーの公式

文章中に埋めこまれた $e^{i\theta}=\cos\theta+i\sin\theta$ オイラーの公式。

\end{enumerate}

\includegraphics{latex.png}

\end{document}


\section{レイアウト}
\layout

% \iffalse
%</doc>
% \fi
%
% \Finale
\endinput
%
% \CheckSum{0}
%
% \CharacterTable
%  {Upper-case    \A\B\C\D\E\F\G\H\I\J\K\L\M\N\O\P\Q\R\S\T\U\V\W\X\Y\Z
%   Lower-case    \a\b\c\d\e\f\g\h\i\j\k\l\m\n\o\p\q\r\s\t\u\v\w\x\y\z
%   Digits        \0\1\2\3\4\5\6\7\8\9
%   Exclamation   \!     Double quote  \"     Hash (number) \#
%   Dollar        \$     Percent       \%     Ampersand     \&
%   Acute accent  \'     Left paren    \(     Right paren   \)
%   Asterisk      \*     Plus          \+     Comma         \,
%   Minus         \-     Point         \.     Solidus       \/
%   Colon         \:     Semicolon     \;     Less than     \<
%   Equals        \=     Greater than  \>     Question mark \?
%   Commercial at \@     Left bracket  \[     Backslash     \\
%   Right bracket \]     Circumflex    \^     Underscore    \_
%   Grave accent  \`     Left brace    \{     Vertical bar  \|
%   Right brace   \}     Tilde         \~}
%
%
